\chapter{Succinct Data Structures}\label{ch22:succinct-data-structures}

\section{Introduction}\label{sec22:introduction}

A data structure is \textbf{succinct} if it stores data in space close to the information-theoretic lower bound while supporting queries efficiently. Given a set of $n$ elements drawn from a universe $\{0, 1, \ldots, m-1\}$, the minimum number of bits needed to represent the set is $n \log(m/n)$ bits (this follows from counting the number of such sets and taking the logarithm). A succinct representation uses this lower bound plus $o(n)$ extra bits, while still supporting the required operations.

\begin{itemize}
\item \textbf{Implicit:} stores the set in exactly the information-theoretic minimum bits (no extra space). Usually cannot support any non-trivial queries.
\item \textbf{Succinct:} uses $H + o(n)$ bits where $H$ is the information-theoretic minimum. Supports queries in O(1) or O(log n) time.
\item \textbf{Compressed:} uses O(n) bits (or $O(H)$ bits), but does not necessarily support efficient queries. A compressed data structure may use more space than a succinct one if $H \ll n$.
\end{itemize}

A concrete comparison: storing an array of $n$ integers from $\{0, \ldots, 2^{32}-1\}$ naively costs $32n$ bits. If the array is a permutation, the information-theoretic minimum is $n \log n$ bits (about $20n$ bits for $n = 10^6$). A succinct representation stores the permutation in $n \log n + o(n)$ bits with O(1) access time.

\textbf{Applications:}
\begin{itemize}
\item Full-text indexes: succinct suffix arrays and FM-indexes index terabyte-scale text in a fraction of the original size.
\item Web search: search engines use succinct data structures to index billions of pages within memory limits.
\item Bioinformatics: genome sequences (billions of base pairs) are indexed using FM-indexes built on rank/select.
\item Compressed databases: SQLite and other embedded databases use succinct structures for query execution on compressed data.
\end{itemize}

\section{Rank and Select on Bitvectors}\label{sec22:rank-select}

Two fundamental operations on bitvectors are \textbf{rank} and \textbf{select}. They form the building blocks of almost all succinct data structures.

\begin{itemize}
\item \textbf{rank(b, i):} the number of 1-bits in the first $i$ positions of bitvector $b$. Formally, $\text{rank}_b(i) = |\{j < i : b[j] = 1\}|$.
\item \textbf{select(b, i):} the position of the $i$-th 1-bit in bitvector $b$. Formally, $\text{select}_b(i) = \min\{j : \text{rank}_b(j+1) = i + 1\}$.
\end{itemize}

Both operations can be supported in O(1) time with a small amount of extra space. The key idea is \textbf{block decomposition}: partition the bitvector into blocks of $\log^2 n$ bits each. For each block, store the rank at the block boundary. Within each block, further divide into superblocks of $\log n$ bits and store the rank at superblock boundaries. The total extra space is $O(n / \log n)$ bits, and answering a rank query requires only a constant number of table lookups.

Select can be answered in O(1) time using a similar decomposition combined with a precomputed table, or in O(log n) time using binary search with rank queries. A practical approach uses interpolation search within blocks since the density of 1-bits is known.

\begin{lstlisting}[language=C++]
class succinct_bitvector {
public:
    succinct_bitvector(const std::vector<bool>& bits)
        : n_(bits.size()),
          blocks_((n_ + 63) / 64),
          bits_(blocks_, 0),
          block_rank_(blocks_, 0)
    {
        for (size_t i = 0; i < n_; ++i)
            if (bits[i]) bits_[i / 64] |= 1ULL << (i % 64);

        uint64_t running = 0;
        for (size_t i = 0; i < blocks_; ++i) {
            block_rank_[i] = running;
            running += __builtin_popcountll(bits_[i]);
        }
    }

    uint64_t rank(uint64_t i) const {
        uint64_t block = i / 64;
        uint64_t pos = i % 64;
        uint64_t r = block > 0 ? block_rank_[block - 1] : 0;
        uint64_t mask = (pos == 64) ? ~0ULL : (1ULL << pos) - 1;
        r += __builtin_popcountll(bits_[block] & mask);
        return r;
    }

    uint64_t select(uint64_t k) const {
        uint64_t lo = 0, hi = n_;
        while (lo < hi) {
            uint64_t mid = lo + (hi - lo) / 2;
            if (rank(mid) <= k) lo = mid + 1;
            else hi = mid;
        }
        return lo;
    }

    uint64_t operator[](uint64_t i) const {
        return (bits_[i / 64] >> (i % 64)) & 1;
    }

    uint64_t size() const { return n_; }

private:
    uint64_t n_, blocks_;
    std::vector<uint64_t> bits_;
    std::vector<uint64_t> block_rank_;
};
\end{lstlisting}

\textbf{Applications of rank and select:}
\begin{itemize}
\item Converting between binary and unary representations: rank converts a bitvector to an array of integers (positions of 1-bits); select does the inverse.
\item Counting occurrences in a range: $\text{rank}(r) - \text{rank}(l)$ counts 1-bits in $[l, r)$.
\item Finding the $k$-th occurrence of a symbol in a string (used by wavelet trees and FM-indexes).
\end{itemize}

\section{Wavelet Tree}\label{sec22:wavelet-tree}

The \textbf{wavelet tree} (Grossi, Gupta, and Vitter, 2003) is a multi-level binary partition of an alphabet that supports rank, select, and range count queries in $O(\log \sigma)$ time, where $\sigma$ is the alphabet size. Each node of the tree stores a bitvector marking which elements of the current sub-alphabet go left (lower half) versus right (upper half).

\subsection{Structure}

For the string \texttt{abracadabra} with alphabet $\{a, b, c, d, r\}$:

\begin{lstlisting}[language=Java]
Wavelet tree for "abracadabra":

              [a b c d r]
              /         \
          [a b]         [c d r]
          / \           / \
        [a] [b]      [c] [d r]
                            / \
                          [d] [r]

Root bitvector (left=0, right=1):
  a b r a c a d a b r a
  0 0 1 0 1 0 1 0 0 1 0   (a,b go left; r,c,d go right)

Left child bitvector (a vs b):
  a b . a . a . a b . a
  0 1 . 0 . 0 . 0 1 . 0   (a=0, b=1)

Right child bitvector (c vs d,r):
  . . r . c . d . . r .
  . . 1 . 0 . 0 . . 1 .   (c=0, d,r go right)

Far-right child (d vs r):
  . . r . . . d . . r .
  . . 1 . . . 0 . . 1 .   (d=0, r=1)
\end{lstlisting}

\subsection{Operations}

\textbf{access(i):} starting at the root, read bitvector[pos]. If 0, go left; if 1, go right. The symbol is found at the leaf. Time: $O(\log \sigma)$.

\textbf{rank(c, i):} at the root, determine if $c$ goes left or right. If $c$ goes left, replace $i$ with rank0(i) and recurse left. If $c$ goes right, replace $i$ with rank1(i) and recurse right. Time: $O(\log \sigma)$.

\textbf{select(c, i):} starting at the leaf for $c$, walk upward. At each parent, if $c$ came from the left child, find the position in the parent bitvector whose rank0 equals $i$. If from the right, use rank1. Time: $O(\log \sigma)$.

\textbf{range count(l, r, c\_lo, c\_hi):} count occurrences of symbols in $[c\_lo, c\_hi]$ within $[l, r)$ of the text. This is answered by recursing into the wavelet tree, narrowing the alphabet range at each level. Time: $O(\log \sigma)$.

\begin{lstlisting}[language=C++]
struct wavelet_node {
    std::vector<bool> bits;
    std::vector<uint64_t> block_rank;
    std::unique_ptr<wavelet_node> left, right;
    char lo, hi;

    uint64_t rank0(uint64_t i) const {
        uint64_t r = 0;
        for (uint64_t j = 0; j < i; ++j)
            if (!bits[j]) ++r;
        return r;
    }

    uint64_t rank1(uint64_t i) const {
        return i - rank0(i);
    }
};

class wavelet_tree {
public:
    wavelet_tree(const std::string& s) {
        std::set<char> alpha(s.begin(), s.end());
        root_ = build(s, alpha.begin(), alpha.end());
    }

    char access(uint64_t i) const {
        return access(root_.get(), i);
    }

    uint64_t rank(char c, uint64_t i) const {
        return rank(root_.get(), c, i);
    }

    uint64_t select(char c, uint64_t k) const {
        return select(root_.get(), c, k);
    }

    uint64_t range_count(uint64_t l, uint64_t r,
                         char lo, char hi) const {
        return range_count(root_.get(), l, r, lo, hi);
    }

private:
    std::unique_ptr<wavelet_node> root_;

    std::unique_ptr<wavelet_node> build(
        const std::string& s,
        std::set<char>::iterator lo,
        std::set<char>::iterator hi)
    {
        if (lo == hi) return nullptr;
        if (std::next(lo) == hi) {
            auto node = std::make_unique<wavelet_node>();
            node->lo = node->hi = *lo;
            node->bits.assign(s.size(), false);
            return node;
        }
        auto mid = lo;
        std::advance(mid, std::distance(lo, hi) / 2);
        auto node = std::make_unique<wavelet_node>();
        node->lo = *lo;
        node->hi = *std::prev(hi);

        std::string left_str, right_str;
        for (char c : s) {
            bool go_left = (c < *mid);
            node->bits.push_back(go_left);
            if (go_left) left_str += c;
            else right_str += c;
        }
        node->left = build(left_str, lo, mid);
        node->right = build(right_str, mid, hi);
        return node;
    }

    char access(const wavelet_node* node, uint64_t i) const {
        if (!node) return node->lo;
        if (node->left == nullptr && node->right == nullptr)
            return node->lo;
        if (!node->bits[i]) {
            uint64_t new_i = node->rank0(i);
            return access(node->left.get(), new_i);
        } else {
            uint64_t new_i = node->rank1(i);
            return access(node->right.get(), new_i);
        }
    }

    uint64_t rank(const wavelet_node* node, char c,
                  uint64_t i) const {
        if (!node) return 0;
        if (node->left == nullptr && node->right == nullptr)
            return i;
        if (c <= node->left->hi) {
            uint64_t new_i = node->rank0(i);
            return rank(node->left.get(), c, new_i);
        } else {
            uint64_t new_i = node->rank1(i);
            return rank(node->right.get(), c, new_i);
        }
    }

    uint64_t select(const wavelet_node* node, char c,
                    uint64_t k) const {
        if (!node) return k;
        if (node->left == nullptr && node->right == nullptr)
            return k;
        uint64_t child_sel;
        if (c <= node->left->hi)
            child_sel = select(node->left.get(), c, k);
        else
            child_sel = select(node->right.get(), c, k);
        return node->bits[child_sel] ? node->rank1(child_sel + 1)
                                     : node->rank0(child_sel + 1) - 1;
    }

    uint64_t range_count(const wavelet_node* node,
                         uint64_t l, uint64_t r,
                         char lo, char hi) const {
        if (!node || l >= r) return 0;
        if (hi < node->lo || lo > node->hi) return 0;
        if (lo <= node->lo && node->hi <= hi)
            return r - l;
        uint64_t l0 = node->rank0(l), r0 = node->rank0(r);
        uint64_t l1 = node->rank1(l), r1 = node->rank1(r);
        uint64_t total = 0;
        if (lo <= node->left->hi)
            total += range_count(node->left.get(), l0, r0, lo, hi);
        if (hi >= node->right->lo)
            total += range_count(node->right.get(), l1, r1, lo, hi);
        return total;
    }
};
\end{lstlisting}

\textbf{Complexity:} each operation traverses at most $\log \sigma$ levels, doing O(1) work at each. Space: $O(n \log \sigma)$ bits total across all bitvectors (the sum of bitvector lengths at each level equals $n$). With block decomposition on each bitvector, the overhead is $O(n \log \sigma / \log n)$ extra bits.

\textbf{Applications:} range median queries, range quantile queries, string indexing, frequency counting, and as a component of the FM-index.

\section{Burrows-Wheeler Transform and FM-Index}\label{sec22:bwt-fmindex}

\subsection{Burrows-Wheeler Transform}

The \textbf{Burrows-Wheeler Transform} (BWT) permutes a string by sorting all of its rotations lexicographically. The transformed string is the last column of the sorted rotation matrix. The BWT is reversible: given the transformed string and the original string's starting position in the sorted matrix, the original string can be recovered.

\textbf{Example:} for the string \texttt{banana\$}:

\begin{lstlisting}[language=Java]
  Sorted rotations of "banana$":

  $ a n a n a      Last     First
  a $ a n a n      column   column
  a n $ a n a      ------   ------
  a n a $ a n   -> a n n $  $ a a a
  a n a n $ a      n a a a  a n n a
  a n a n a $      a a n n  n a $ a
  n a $ a n a      n $ a a  n a a $
  n a n a $ a                  $ a n a

  BWT("banana$") = "annb$aa"
  (concatenating last column from top to bottom)
\end{lstlisting}

The BWT does not compress by itself, but it groups similar characters together, making the output highly compressible by run-length encoding or other methods.

\subsection{FM-Index}

The \textbf{FM-Index} (Ferragina and Manzini, 2000) is a succinct full-text index that uses the BWT together with rank on auxiliary bitvectors. It supports substring search in time proportional to the pattern length, using O(n) bits of space.

The FM-Index maintains, for each character $c$, a bitvector $B_c$ of length $n+1$ where $B_c[i] = 1$ if the $i$-th character of the BWT is $c$. A C-array stores the number of characters lexicographically smaller than $c$ in the original text.

\textbf{LF-mapping:} the key property is that the $i$-th occurrence of character $c$ in the BWT corresponds to the $i$-th occurrence of $c$ in the original text. This mapping, called the LF-mapping, is:
\[
\text{LF}(i) = C[\text{BWT}[i]] + \text{rank}_{B_c}(i)
\]
where $C[c]$ counts characters lexicographically smaller than $c$.

\textbf{Backward search:} to find a pattern $P = p_1 p_2 \ldots p_m$, start with an interval $[lo, hi) = [0, n)$ and extend backward. For each character $p_j$ from right to left:
\[
lo = C[p_j] + \text{rank}_{B_{p_j}}(lo), \quad
hi = C[p_j] + \text{rank}_{B_{p_j}}(hi)
\]
If $lo < hi$ after processing all characters, the pattern occurs $hi - lo$ times in the text. Time: O(m) with precomputed rank, or O(m log n) with binary-search rank.

\begin{lstlisting}[language=C++]
class fm_index {
public:
    fm_index(const std::string& text) : n_(text.size()) {
        build_bwt(text);
    }

    int count(const std::string& pattern) const {
        int lo = 0, hi = n_;
        for (int i = (int)pattern.size() - 1; i >= 0; --i) {
            char c = pattern[i];
            lo = C_[c] + rank(c, lo);
            hi = C_[c] + rank(c, hi);
            if (lo >= hi) return 0;
        }
        return hi - lo;
    }

    std::vector<int> find(const std::string& pattern) const {
        int lo = 0, hi = n_;
        for (int i = (int)pattern.size() - 1; i >= 0; --i) {
            char c = pattern[i];
            lo = C_[c] + rank(c, lo);
            hi = C_[c] + rank(c, hi);
            if (lo >= hi) return {};
        }
        std::vector<int> occ;
        for (int i = lo; i < hi; ++i)
            occ.push_back(locate(i));
        return occ;
    }

private:
    int n_;
    std::vector<char> bwt_;
    std::vector<int> C_;
    std::vector<std::vector<int>> occ_;

    void build_bwt(const std::string& text) {
        n_ = text.size();
        std::vector<int> sa(n_);
        std::iota(sa.begin(), sa.end(), 0);
        std::sort(sa.begin(), sa.end(), [&](int a, int b) {
            return text.compare(a, n_ - a, text, b, n_ - b) < 0;
        });

        bwt_.resize(n_);
        for (int i = 0; i < n_; ++i)
            bwt_[i] = text[(sa[i] + n_ - 1) % n_];

        C_.assign(256, 0);
        for (char c : bwt_) ++C_[c];
        int sum = 0;
        for (int i = 0; i < 256; ++i) {
            int tmp = C_[i];
            C_[i] = sum;
            sum += tmp;
        }

        occ_.assign(256, std::vector<int>(n_ + 1, 0));
        for (int i = 0; i < n_; ++i) {
            for (int c = 0; c < 256; ++c)
                occ_[c][i + 1] = occ_[c][i];
            ++occ_[(int)bwt_[i]][i + 1];
        }
    }

    int rank(char c, int i) const {
        return occ_[(int)c][i];
    }

    int locate(int pos) const {
        int steps = 0;
        while (steps < n_) {
            char c = bwt_[pos];
            pos = C_[c] + occ_[(int)c][pos];
            ++steps;
        }
        return (pos + 1) % n_;
    }
};
\end{lstlisting}

\textbf{Applications:}
\begin{itemize}
\item Bioinformatics: FM-indexes index entire human genomes (3 billion base pairs) in a few gigabytes, enabling fast substring search for read alignment.
\item Compression: the BWT is the core of bzip2; FM-indexes combine compression with search.
\item Full-text search: the FM-index is the basis of the succinct index in many text search libraries.
\end{itemize}

\section{Succinct Trees}\label{sec22:succinct-trees}

A tree with $n$ nodes uses $O(n \log n)$ bits with pointer-based representations (each pointer is $\log n$ bits). A succinct representation stores the tree in $2n + o(n)$ bits while supporting navigation operations in O(1) time.

\subsection{Balanced Parentheses}

A rooted ordered tree can be encoded as a balanced parentheses sequence: each node is represented by an opening parenthesis \texttt{(} (preorder visit) and a closing parenthesis \texttt{)} (postorder visit). For a tree with $n$ nodes, the sequence has exactly $2n$ parentheses.

\begin{lstlisting}[language=Java]
Tree:           Encoding:
      1         ( ( ) ( ( ) ( ) )
     / \        
    2   3       Node 1: ( ... )
   /   / \      
  4   5   6     Node 2: ( )        -- leaf
      |   |     
      7   8     Node 3: ( ( ) ( ) )
\end{lstlisting}

\textbf{Navigation with rank and select:}
\begin{itemize}
\item \textbf{parent(u):} find the matching open-parenthesis before position $u$ using rank and select.
\item \textbf{first\_child(u):} if $b[u+1] = 1$ (opening), then $u+1$ is the first child. Otherwise, $u$ has no children.
\item \textbf{next\_sibling(u):} find the matching close-parenthesis of $u$, then the next position is the next sibling.
\item \textbf{subtree\_size(u):} the number of pairs within the matched parentheses is the subtree size.
\end{itemize}

All these operations use O(1) rank/select queries, giving O(1) navigation time.

\subsection{LOUDS}

\textbf{Level-Order Unary Degree Sequence} (LOUDS) encodes a tree by listing each node's degree in unary (the degree $d$ is written as $d$ ones followed by a zero). For a tree with $n$ nodes and total degree $\sum d = n-1$, the sequence has $2n$ bits.

\begin{lstlisting}[language=C++]
class succinct_tree_bp {
public:
    succinct_tree_bp(const std::vector<std::vector<int>>& adj,
                     int root) {
        n_ = 0;
        encode(adj, root);
        build_rank();
    }

    int parent(int u) const {
        int match = find_matching(u);
        return select0(rank0(match)) / 2;
    }

    int first_child(int u) const {
        if (parens_[u + 1] == 1) return u + 1;
        return -1;
    }

    int subtree_size(int u) const {
        int match = find_matching(u);
        return (match - u + 1) / 2;
    }

    int depth() const { return max_depth_; }

    size_t memory_bits() const {
        return parens_.size() + block_rank_.size() * 64;
    }

private:
    int n_;
    std::vector<bool> parens_;
    std::vector<uint64_t> block_rank_;
    int max_depth_ = 0;

    void encode(const std::vector<std::vector<int>>& adj,
                int u) {
        parens_.push_back(true);
        ++n_;
        int depth = 0;
        for (int v : adj[u]) {
            encode(adj, v);
            ++depth;
        }
        parens_.push_back(false);
        max_depth_ = std::max(max_depth_, depth);
    }

    void build_rank() {
        uint64_t blocks = (parens_.size() + 63) / 64;
        block_rank_.resize(blocks, 0);
        uint64_t running = 0;
        for (uint64_t i = 0; i < blocks; ++i) {
            block_rank_[i] = running;
            uint64_t start = i * 64;
            uint64_t end = std::min(start + 64, parens_.size());
            for (uint64_t j = start; j < end; ++j)
                if (parens_[j]) ++running;
        }
    }

    uint64_t rank0(uint64_t i) const {
        return i - rank1(i);
    }

    uint64_t rank1(uint64_t i) const {
        uint64_t block = i / 64;
        uint64_t pos = i % 64;
        uint64_t r = block > 0 ? block_rank_[block - 1] : 0;
        uint64_t start = block * 64;
        uint64_t mask = (pos == 64) ? ~0ULL : (1ULL << pos) - 1;
        uint64_t word = 0;
        for (uint64_t j = 0; j < 64 && start + j < parens_.size(); ++j)
            if (parens_[start + j]) word |= 1ULL << j;
        r += __builtin_popcountll(word & mask);
        return r;
    }

    uint64_t select0(uint64_t k) const {
        uint64_t lo = 0, hi = parens_.size();
        while (lo < hi) {
            uint64_t mid = lo + (hi - lo) / 2;
            if (rank0(mid) <= k) lo = mid + 1;
            else hi = mid;
        }
        return lo;
    }

    int find_matching(int pos) const {
        int depth = 0;
        for (int i = pos; i < (int)parens_.size(); ++i) {
            if (parens_[i]) ++depth;
            else {
                --depth;
                if (depth == 0) return i;
            }
        }
        return -1;
    }
};
\end{lstlisting}

\textbf{Space comparison:} a pointer-based binary tree uses $O(n \log n)$ bits. The balanced-parentheses and LOUDS representations use exactly $2n$ bits ($2n + o(n)$ with auxiliary structures). For $n = 10^6$, the pointer-based tree uses about 24 MB; the succinct representation uses about 250 KB.

\section{Compressed Suffix Arrays}\label{sec22:compressed-suffix-arrays}

A \textbf{suffix array} stores the lexicographic order of all $n$ suffixes of a string of length $n$. The naive representation uses $n \log n$ bits (one integer per suffix). A \textbf{compressed suffix array} stores this information in O(n) bits while supporting the same queries.

The FM-Index is one form of compressed suffix array: it encodes the suffix array implicitly through the BWT. Other forms include:
\begin{itemize}
\item \textbf{SSA} (Succinct Suffix Array): stores the suffix array using a wavelet tree on the BWT, supporting rank, select, and pattern matching.
\item \textbf{PSA} (Compressed Suffix Array): stores the suffix array as a sequence of differences and compresses using wavelet trees or entropy-bounded integers.
\item \textbf{LZ-indexes:} based on Lempel-Ziv decomposition rather than suffix arrays, offering different space-time tradeoffs.
\end{itemize}

All these structures achieve O(n) bits of space and O(m) or O(m log n) pattern matching time, where $m$ is the pattern length.

\section{Applications and Industry Use}\label{sec22:applications}

\begin{itemize}
\item \textbf{Search engines:} Google and Bing use succinct indexes for web-scale full-text search. The FM-index variant indexes trillions of tokens within memory budgets that would be impossible with naive structures.

\item \textbf{Genome browsers:} tools like BWA and Bowtie use FM-indexes to align short DNA reads to reference genomes. The human genome (3 billion base pairs) is indexed in about 4 GB using an FM-index.

\item \textbf{Compressed databases:} SQLite and LevelDB use succinct bitvectors and wavelet trees for compressed column storage and range queries without decompressing the entire dataset.

\item \textbf{In-memory caches:} Redis modules and Memcached plugins use succinct data structures to increase cache capacity without adding memory hardware.

\item \textbf{Log analysis:} succinct wavelet trees support frequency histograms and range quantiles on compressed log data, enabling analytics without full decompression.
\end{itemize}

\section{Common Bugs and Pitfalls}\label{sec22:pitfalls}

\begin{itemize}
\item \textbf{Off-by-one in rank/select:} rank and select may be defined as 0-indexed or 1-indexed depending on the convention. Mixing conventions within the same code produces silent errors. Pick one convention and use it consistently.

\item \textbf{select boundary conditions:} select(k) where $k$ exceeds the number of 1-bits is undefined. Always check that $k < \text{rank}(n)$ before calling select.

\item \textbf{Wavelet tree alphabet partitioning:} if the alphabet is not partitioned correctly at each level, rank and select return wrong values. The split point must divide the current sub-alphabet roughly in half to maintain $O(\log \sigma)$ depth.

\item \textbf{BWT construction with sentinel:} the BWT requires a unique sentinel character (e.g., \texttt{\$}) at the end of the string. Omitting the sentinel causes the first and last rotations to collide in the sorted matrix.

\item \textbf{Block size tuning:} for bitvector rank, block size $\log^2 n$ is theoretically optimal but may be too large for small bitvectors. For bitvectors under 1024 bits, a simple O(n) scan for rank is faster than block decomposition.

\item \textbf{Memory alignment:} succinct data structures pack bits tightly, which can cause unaligned memory access on some architectures. Use compiler intrinsics (\texttt{\_\_builtin\_popcountll}) and ensure arrays are allocated on aligned boundaries.
\end{itemize}

\section{Complexity Summary}\label{sec22:complexity-summary}

\begin{center}
\begin{tabular}{lccc}
\toprule
Operation & Time & Space & Structure \\
\midrule
Bitvector rank(b, i) & O(1) & O(n / log n) extra & Block decomposition \\
Bitvector select(b, i) & O(1) & O(n / log n) extra & Block decomposition + table \\
Wavelet tree access(i) & O(log $\sigma$) & O(n log $\sigma$) bits & Wavelet tree \\
Wavelet tree rank(c, i) & O(log $\sigma$) & O(n log $\sigma$) bits & Wavelet tree \\
Wavelet tree select(c, k) & O(log $\sigma$) & O(n log $\sigma$) bits & Wavelet tree \\
Wavelet tree range count & O(log $\sigma$) & O(n log $\sigma$) bits & Wavelet tree \\
FM-index pattern count & O(m) & O(n) bits & BWT + rank bitvectors \\
FM-index locate & O(m + occ) & O(n) bits & BWT + rank bitvectors \\
Succinct tree parent & O(1) & 2n + o(n) bits & Balanced parentheses \\
Succinct tree first\_child & O(1) & 2n + o(n) bits & Balanced parentheses \\
Succinct tree subtree\_size & O(1) & 2n + o(n) bits & Balanced parentheses \\
\bottomrule
\end{tabular}
\end{center}

\section{Exercises}\label{sec22:exercises}

\begin{enumerate}
\item Implement rank and select on a bitvector of 1 million random bits. Measure the space overhead (in bits per element) and verify that rank answers in O(1) time empirically by timing 100{,}000 random queries.

\item Build a wavelet tree for the string \texttt{mississippi}. Verify that rank(\texttt{s}, 7) returns the correct count and that select(\texttt{s}, 3) returns the position of the third \texttt{s}. Trace the algorithm by hand and confirm the result matches.

\item Construct the BWT of the string \texttt{cacao\$} by hand: list all rotations, sort them, and read the last column. Then verify that you can recover the original string using the LF-mapping.

\item Implement an FM-index for a small text (at most 1000 characters). Test pattern matching on \texttt{abracadabra} and verify that the count for \texttt{abra} returns 2.

\item Compare the space of a pointer-based binary tree with $n = 10^5$ nodes to the balanced-parentheses succinct representation. Print both sizes in bytes.

\item Prove that a bitvector of length $n$ with $k$ set bits can be represented in $n H_0 + o(n)$ bits, where $H_0 = -\frac{k}{n}\log_2\frac{k}{n} - \frac{n-k}{n}\log_2\frac{n-k}{n}$ is the zeroth-order entropy. Hint: this is the number of distinct bitvectors with exactly $k$ ones.

\item Implement a succinct bitvector with block size $\log^2 n$ and compare its rank performance with a naive O(n) scan for bitvectors of size $10^4, 10^5, 10^6$. At what size does block decomposition become faster?

\item Given a wavelet tree for a DNA string (alphabet $\{A, C, G, T\}$), implement a function that returns the frequency of each base in a given range $[l, r)$. Analyse the complexity and compare with a prefix-sum approach.
\end{enumerate}

\section{References and Further Reading}\label{sec22:references}

\begin{itemize}
\item Jacobson, G. (1989). ``Succinct Static Data Structures.'' PhD thesis, Carnegie Mellon University. The foundational work on succinct data structures, introducing balanced parentheses and LOUDS encodings for trees.

\item Grossi, R., Gupta, A., and Vitter, J. S. (2003). ``High-Order Entropy-Compressed Text Indexes.'' Proceedings of the 14th Annual ACM-SIAM Symposium on Discrete Algorithms (SODA). Introduces the wavelet tree and its applications to text indexing.

\item Ferragina, P. and Manzini, G. (2005). ``An Alphabet-Friendly Succinct Index.'' Journal of the ACM, 52(3). Introduces the FM-index, showing how the BWT combined with rank/select gives O(n) bit full-text indexing.

\item Sadakane, K. (2007). ``Succinct Representations of Large Graphs and Strings.'' Proceedings of the 18th Annual ACM-SIAM Symposium on Discrete Algorithms (SODA). Extends succinct representations to trees, planar graphs, and strings.

\item Navarro, G. and M\"akinen, V. (2011). ``Compressed Full-Text Indexes.'' ACM Computing Surveys, 43(4). A comprehensive survey of compressed indexes including wavelet trees, FM-indexes, and LZ-based indexes.

\item Munro, J. I. and Raman, V. (1997). ``Succinct Representation of Permutations and Functions.'' Proceedings of the 23rd International Colloquium on Automata, Languages, and Programming (ICALP). Shows how permutations and functions can be represented in $n \log n + O(n)$ bits with O(1) query time.

\item Clark, D. (1996). ``Compact Pat Trees.'' PhD thesis, University of Waterloo. Describes succinct representations of tries and pat trees using rank/select.
\end{itemize}
